% Integration fragment: uses only standard amsmath/booktabs already in the paper.
% Proposed central result; no empirical outcome is asserted here.
\subsection{What a task can identify without a calibrated neural response curve}
\label{sec:monotone_identification}

A successful linear-response model comparison does not establish that the latent
computational quantities are distinguishable independently of the assumed neural
encoding. We therefore first characterize identification at the level of task
support. Let $f_i$ be a scalar candidate's prediction in condition $i$, where a
condition includes the relevant outcome, confidence, and history. The latent
parameters are held fixed or constrained independently by behavior. Allow a
single, candidate-specific, strictly increasing response function $h$, shared
across conditions, with unrestricted offset and dynamic range. Its possible
premeasurement mean vectors form
\begin{equation}
\mathcal C_f=\{(h(f_1),\ldots,h(f_n)):h\text{ strictly increasing}\}.
\end{equation}

\paragraph{Finite-support equivalence.}
Two candidates have identical premeasurement mean families if and only if their
predictions have the same weak order, including the same ties:
\begin{equation}
\mathcal C_f=\mathcal C_g
\quad\Longleftrightarrow\quad
\operatorname{sign}(f_i-f_j)=\operatorname{sign}(g_i-g_j)
\quad\text{for every }i,j.
\label{eq:weak_order_identification}
\end{equation}
Indeed, an increasing response preserves every strict inequality and equality.
Conversely, identical weak orders define an increasing bijection between the
observed candidate levels; interpolation and composition with the response
function give the same set of mean vectors. If the weak orders differ, the two
families are disjoint before measurement, and a two-condition witness exists:
either the candidates reverse their ordering, or one predicts equality where
the other predicts a difference. Thus a matched-level contrast can identify a
distinction that an arbitrarily large number of repetitions on an equivalent
support cannot. Imperfect Pearson correlation between candidate regressors is
not sufficient: a nonlinear increasing reparameterization can leave their
entire encoded mean families identical.

\paragraph{A five-candidate equivalence on fixed-confidence binary tasks.}
For a binary transition with observed-outcome probability $q$ and Dirichlet
concentration $a$, the canonical quantities in Table~\ref{tab:candidates} reduce to
\begin{align}
I&=\sqrt{2}(1-q), & S&=-\log q, & G&=(a+1)^{-1},\nonumber\\
U&=\frac{\sqrt{2}(1-q)}{a+1}, &
K&=-\log q+\psi(aq+1)-\psi(a+1).
\label{eq:canonical_binary_candidates}
\end{align}
Here $K=\mathrm{KL}[\operatorname{Dir}(\boldsymbol\alpha+\mathbf e_j)
\Vert\operatorname{Dir}(\boldsymbol\alpha)]$. For a one-step stable/reset
mixture with fixed hazard $\eta\in(0,1)$ and fixed reset predictive probability
$r_j>0$, the reset posterior is
\begin{equation}
R=\frac{\eta r_j}{(1-\eta)q+\eta r_j}.
\label{eq:restricted_reset_posterior}
\end{equation}
At fixed $a,\eta,r_j$, all five nonconstant quantities $I,S,U,K,R$ are strictly
decreasing in $q$ and hence constitute one increasing-link equivalence class;
$G$ is constant. For the less immediate case,
$\partial K/\partial q=-1/q+a\psi_1(aq+1)<0$, because
$\psi_1(x+1)=\sum_{k\geq1}(x+k)^{-2}<1/x$ for $x>0$.
Consequently, recovering a winner among these five on this support using linear
response curves identifies a winner \emph{within that encoding class}, not a
distinction invariant to the neural response curve.

\paragraph{A minimal separating support.}
Varying confidence independently of probability separates four classes in the
binary, fixed-reset setting. A concrete three-condition support is
$A:(q,a)=(0.8,5)$, $B:(0.1,50)$, and $C:(0.8,50)$:
\begin{center}
\begin{tabular}{llll}
\toprule
Quantity or class & Predicted ordering & Separating feature\\
\midrule
$I,S,R$ & $A=C<B$ & Matched probability\\
$G$ & $B=C<A$ & Matched concentration\\
$U$ & $C<B<A$ & Update magnitude\\
$K$ & $C<A<B$ & Opposite $A$--$B$ ordering to $U$\\
\bottomrule
\end{tabular}
\end{center}
In particular, $(U_A,U_B,U_C)=(0.04714045,0.02495671,0.00554594)$,
whereas $(K_A,K_B,K_C)=(0.02314355,0.08671309,0.00248125)$.
These are deterministic predictions, not fitted or empirical measurements.
The three-condition support is cardinality-minimal for these four classes:
two conditions permit only three weak-order signatures ($<$, $=$, or $>$).
This result concerns support, not sample allocation or power, and assumes that
the intended probability and confidence values are independently established.

The remaining equivalences have explicit boundaries. With three or more
outcomes, $I^2=(1-q)^2+\sum_{k\ne j}p_k^2$ is not generally a function of $q$.
For example, observing category 1 under $(0.5,0.49,0.01)$ versus
$(0.5,0.25,0.25)$ holds surprise and, at matched concentration, $K$ fixed while
changing innovation. Similarly, independently changing reset hazard or reset
predictive probability at matched $q$ breaks
Eq.~\ref{eq:restricted_reset_posterior}'s equivalence to surprise. The restricted
reset result is not a general claim about a history-dependent reset filter.

For completeness, three conditions also suffice to separate all six declared
scalar quantities when three-outcome predictive distributions and reset hazards
can vary independently. With category 1 observed and $r_j=1/3$, take
$A:(\mathbf p,a,\eta)=((0.8,0.1,0.1),5,0.1)$,
$B:((0.1,0.45,0.45),50,0.1)$, and
$C:((0.8,0.19,0.01),50,0.8)$. The respective weak orders are
$A<C<B$ for $I$, $A=C<B$ for $S$, $B=C<A$ for $G$, $C<B<A$ for $U$,
$C<A<B$ for $K$, and $A<B<C$ for the one-step reset posterior. This constructive
example establishes a mathematical support bound, not the practical feasibility
of independently manipulating these latent quantities in an animal.

\paragraph{Measurement imposes a second identification condition.}
Let $H$ be a fixed known linear event-to-sample operator incorporating sensor
convolution, and let $B\boldsymbol\beta$ be unrestricted linear nuisance.
The measured mean family is
$\mathcal M_f=H\mathcal C_f+\operatorname{col}(B)$.
Writing $P$ for projection orthogonal to $\operatorname{col}(B)$, two families
overlap exactly when
\begin{equation}
\exists\,\mathbf z\in\mathcal C_f,\ \mathbf w\in\mathcal C_g:
\quad PH(\mathbf z-\mathbf w)=0.
\label{eq:measured_identification}
\end{equation}
Equal weak orders therefore remain observationally equivalent after every common
sensor/nuisance mapping. Distinct weak orders need not remain distinguishable:
observing only the sum of two conditions erases the distinction between
$(0,1)$ and $(1,0)$. By contrast, a full-column-rank sensor with no nuisance
preserves structural distinctions even if it makes their estimation
ill-conditioned. More generally, $\ker(PH)\subseteq\operatorname{span}(\mathbf1)$
is sufficient to preserve all weak-order distinctions, because a common offset
cannot alter an ordering.

Equation~\ref{eq:measured_identification} gives a reproducible finite diagnostic.
Let $T_f$ and $T_g$ assign each condition to its ordered candidate level, so
$\mathbf z=T_f\mathbf u$ and $\mathbf w=T_g\mathbf v$. Test
$PH(T_f\mathbf u-T_g\mathbf v)=0$ with adjacent level gaps at least one.
Homogeneity makes these constraints equivalent to strictly positive gaps for
the unrestricted-link model. To avoid amplification of projection-cancellation
roundoff, the implementation solves the equivalent direct equation
$H(T_f\mathbf u-T_g\mathbf v)+B\boldsymbol\gamma=0$ with free nuisance
coefficients, one global sensor normalization, and columnwise nuisance
normalization. It checks numerical feasibility and residuals; it is not an exact-arithmetic
certificate. An unknown fitted kernel or hypothesis-specific latent parameters
requires an additional union over those parameters and can create further
overlap.

\paragraph{A witness is not a power guarantee.}
Unrestricted increasing links allow the dynamic range to approach zero:
$h_\epsilon(x)=b+\epsilon x$ remains increasing for every $\epsilon>0$.
For shared Gaussian noise $\sigma^2 I$, the divergence between the corresponding
measured distributions is
\begin{equation}
D_{\mathrm{KL}}=
\frac{\epsilon^2}{2\sigma^2}\|H(\mathbf f-\mathbf g)\|_2^2
\longrightarrow0.
\end{equation}
Thus no fixed finite sample size gives uniformly positive discrimination power
above test size over unrestricted increasing links. A robust power or design
claim needs a stated minimum effect scale in addition to a separating support
and an adequate measurement operator.

Finally, the increasing orientation is an assumption: allowing unknown sign
also identifies an ordering with its reverse, and allowing arbitrary
nondecreasing links gives all candidates a constant-response submodel.
Experimental near-matches are not mathematical ties; matching uncertainty must
be propagated or an equivalence tolerance specified independently.
The supplied support audit therefore reports numerical rank/tie witnesses and
near-ties separately, rather than treating an insignificant effect as equality.
